\section{Ethical Considerations}\label{sec:ethics}

This study was designed around a binding constraint: uncontacted peoples cannot consent to be studied, and any approach --- physical or informational --- that increases contact risk is unacceptable. Brazil's no-contact policy (in force since 1987), ILO Convention 169, and UNDRIP frame isolation as a choice that outsiders must respect \citep{pringle2014,ilo169,undrip}. Accordingly:

\begin{enumerate}[leftmargin=1.4em, itemsep=1pt]
\item \textbf{No contact.} The study uses exclusively remote, public, aggregate data. No field work, overflights, interviews, or interactions of any kind are proposed or implied.
\item \textbf{No location disclosure.} Coordinates of isolated-peoples records and occupation areas are used only internally to link records to official territory boundaries. No figure, table, or repository artefact publishes settlement or record locations; all outputs are aggregated to Indigenous-Land boundaries (public policy artefacts) or to regional summaries.
\item \textbf{Respectful framing.} We describe peoples, not ``tribes'', and make no claims about lifeways; the object of measurement is land and threat, not people.
\item \textbf{Data protection.} The raw record layer is excluded from the repository by default.
\end{enumerate}

\section{Conclusions and Future Research}\label{sec:conclusion}

Using only open data, we provide the first territory-level quantitative assessment of forest-loss pressure on Brazil's Indigenous Lands hosting uncontacted peoples. Between 2000 and 2023 these 81 lands lost 1.01\% of their year-2000 tree cover internally --- 3.6 times less than their immediate surroundings, 5.3 times less than wider regional buffers, and 2.5 times less than matched non-isolated Indigenous Lands --- confirming their exceptional buffering role. But this buffer is under concentrated assault: internal loss tripled after 2011, pressure concentrates sharply (top quintile: 63.5\% of aggregate ring loss) in documented crisis sites, and the state's protection apparatus was audited in 2024 as under-staffed and leaving twelve records wholly unprotected.

Future work should pursue four directions: (i) causal designs exploiting restriction-of-use ordinances and demarcation timing (staggered difference-in-differences over the territory panel) to separate protection from remoteness; (ii) extension to Peru, Colombia, and Ecuador via RAISG's isolated-peoples layers, testing external validity beyond Brazil; (iii) integration of fire and degradation products (active-fire archives, canopy-height change) to capture sub-canopy pressure invisible to gross-clearing products; and (iv) a systematic, ethics-preserving assessment of the \emph{protection-coverage gap} --- records without any territorial instrument --- by joining the audited record registry to protection-layer geometry.

\section*{Declarations}
\textbf{Funding.} None. \textbf{Conflicts of interest.} None declared. \textbf{Data and code availability.} All data are public (Table~\ref{tab:data}); all code, configuration, result files, and the audit trail are available in the open repository accompanying this preprint. \textbf{AI-assistance disclosure.} AI tools were used for code development and manuscript drafting under the author's direction; the author verified all results, retains full responsibility for the work, and discloses this assistance in line with venue policy. \textbf{Ethics.} No human subjects, no animals, no field work; see Section~\ref{sec:ethics}.

